% Only notation introduced in this reading is registered here.
% Starred uses mark the defining occurrence; no manually maintained page lists.
\providecommand\NullproductsSymbolLegend{1}
% Opt-in, LuaLaTeX-only page-local notation legend. Load in a reading edition.
\RequirePackage{iftex}
\RequireLuaTeX
\RequirePackage{luacode}
\RequirePackage{hyperref}
\RequirePackage{xcolor}
\makeatletter
\newif\ifSymbolLegendEnabled
\newif\ifSymbolLegendCompact
\newif\ifSL@recording
\SL@recordingtrue
\newlength\SymbolLegendWidth
\newlength\SymbolLegendGap
\newlength\SymbolLegendInnerMargin
\newlength\SymbolLegendEdgeMargin
\setlength\SymbolLegendWidth{42mm}
\setlength\SymbolLegendGap{7mm}
\setlength\SymbolLegendInnerMargin{20mm}
\setlength\SymbolLegendEdgeMargin{10mm}
\newbox\SL@box
\newbox\SL@symbolbox
\newbox\SL@entrybox
\newdimen\SL@x
\newdimen\SL@y
\newdimen\SL@available
\def\SLPageEntries{}
\directlua{symbollegend = dofile('tex/impl/symbol-legend.lua')}

% Register once: [prior] means introduced in an earlier volume.
% Otherwise use an introduction marker in the actual defining notation.
% The last argument is an existing local or xr-hyper definition label.
\NewDocumentCommand\SymbolLegendDeclare{O{local} m m m m}{%
  \@ifundefined{SL@symbol@#2}{}{%
    \PackageError{symbol-legend}{Duplicate symbol ID #2}{Use a unique ID.}}%
  \expandafter\gdef\csname SL@symbol@#2\endcsname{#3}%
  \expandafter\gdef\csname SL@meaning@#2\endcsname{#4}%
  \expandafter\gdef\csname SL@label@#2\endcsname{#5}%
  \directlua{
    local mode = '\luaescapestring{#1}'
    if not (mode == 'prior' or mode == 'local') then
      tex.error('symbol-legend: introduction mode must be prior or local')
    end
    symbollegend.declare('\luaescapestring{#2}', mode == 'prior')
  }%
}
% Resolve semantic definition IDs with the manuscript's existing registry.
\NewDocumentCommand\SymbolLegendDeclareID{O{local} m m m m}{%
  \directlua{
    local r = thmlookup and thmlookup.by_id['\luaescapestring{#5}']
    if not r or not (r.env == 'definition') then
      tex.error('symbol-legend: missing definition ID #5')
    else
      token.set_macro('SLResolvedLabel', r.label, 'global')
    end
  }%
  \begingroup
  \edef\SL@next{\endgroup
    \noexpand\SymbolLegendDeclare[#1]{#2}{\unexpanded{#3}}%
      {\unexpanded{#4}}{\SLResolvedLabel}}%
  \SL@next
}
% A local reading introduction has no theorem/definition number. Its label and
% destination are created only when the first introduction marker really ships.
\NewDocumentCommand\SymbolLegendDeclareReading{m m m}{%
  \@ifundefined{SL@symbol@#1}{}{%
    \PackageError{symbol-legend}{Duplicate symbol ID #1}{Use a unique ID.}}%
  \expandafter\gdef\csname SL@symbol@#1\endcsname{#2}%
  \expandafter\gdef\csname SL@meaning@#1\endcsname{#3}%
  \expandafter\gdef\csname SL@label@#1\endcsname{slreading:#1}%
  \expandafter\gdef\csname SL@reading@#1\endcsname{1}%
  \directlua{symbollegend.declare('\luaescapestring{#1}', false, true)}%
}
% Internal shipout callback. The Lua collector inserts the actual PDF anchor
% at the first retained introduction node; the page is expanded here, never at
% box construction. Immediate AUX writing avoids a second delayed shipout.
\newcommand\SLReadingIntroduction[4]{%
  \begingroup
  \protected@edef\SL@readingpage{#2}%
  \expandafter\xdef\csname r@slreading:#1\endcsname{%
    {}{\SL@readingpage}{Einführung}{#3}{}}%
  \if@filesw
    \immediate\write\@auxout{\string\newlabel{slreading:#1}%
      {{}{\SL@readingpage}{Einführung}{#3}{}}}%
  \fi
  \endgroup
}
% Put markers INSIDE the math atom that renders the registered symbol.
% An introduction marker is not a use; the two operations are separate.
\NewDocumentCommand\SymbolLegendUse{m}{%
  \ifSymbolLegendEnabled\ifSL@recording
    \directlua{symbollegend.mark('\luaescapestring{#1}', false)}%
  \fi\fi
}
\NewDocumentCommand\SymbolLegendIntroduce{m}{%
  \ifSymbolLegendEnabled\ifSL@recording
    \directlua{symbollegend.mark('\luaescapestring{#1}', true)}%
  \fi\fi
}
\NewDocumentCommand\SymbolLegendWithoutRecording{m}{%
  \begingroup\SL@recordingfalse #1\endgroup
}
\pdfstringdefDisableCommands{%
  \def\SymbolLegendUse#1{}\def\SymbolLegendIntroduce#1{}%
}

\newcommand\SymbolLegendEntry[1]{%
  \edef\SL@label{\csname SL@label@#1\endcsname}%
  \@ifundefined{r@\SL@label}{%
    \PackageWarning{symbol-legend}{Unresolved definition label \SL@label.
      Build/import the definition and rerun before publishing}}{}%
  \setbox\SL@symbolbox=\hbox{\csname SL@symbol@#1\endcsname}%
  \ifdim\wd\SL@symbolbox>\linewidth
    \PackageError{symbol-legend}{Symbol #1 exceeds the legend width}%
      {Increase SymbolLegendWidth or shorten its display notation.}%
  \fi
  % Link the complete entry, including the symbol and explanation. A single
  % fixed-width box gives PDF readers one generous, non-overlapping tap area.
  \setbox\SL@entrybox=\vbox{%
    \hsize=\linewidth
    \ifSymbolLegendCompact
      \noindent\box\SL@symbolbox\enspace
      \csname SL@meaning@#1\endcsname\par\nobreak
    \else
      \noindent\box\SL@symbolbox\par\nobreak
      \noindent\csname SL@meaning@#1\endcsname\par\nobreak
    \fi
    \noindent
    \ifcsname SL@reading@#1\endcsname
      Einführung, S.~\pageref*{\SL@label}%
    \else
      Def.~\ref*{\SL@label}, S.~\pageref*{\SL@label}%
    \fi\par}%
  \directlua{symbollegend.assert_box_fits(\number\SL@entrybox,
    \number\linewidth, '\luaescapestring{#1}')}%
  \noindent\hyperref[\SL@label]{\box\SL@entrybox}\par
  \ifSymbolLegendCompact\vskip4pt\else\vskip7pt\fi
}
\newcommand\SL@render{%
  \ifx\SLPageEntries\@empty\else
    \begingroup
    \SL@recordingfalse
    \setbox\SL@box=\vbox{%
      % Shipout may interrupt a colored box; never inherit that color state.
      \color{black}%
      \hsize=\SymbolLegendWidth\linewidth=\hsize
      \parindent=0pt\parskip=0pt
      \normalfont\fontsize{9}{11}\selectfont\raggedright
      \textbf{Symbole dieser Seite}\par\vskip6pt
      \SLPageEntries
    }%
    \directlua{symbollegend.assert_box_fits(\number\SL@box,
      \number\SymbolLegendWidth)}%
    \ifdim\dimexpr\ht\SL@box+\dp\SL@box\relax>\textheight
      \PackageError{symbol-legend}{Legend too tall on page \thepage}%
        {Increase the reserved space or shorten the registered meanings.
         No entries may be silently clipped or moved to another page.}%
    \fi
    \SL@x=\dimexpr1in+\hoffset+\oddsidemargin+\textwidth+\SymbolLegendGap\relax
    \if@twoside\ifodd\c@page\else
      \SL@x=\dimexpr1in+\hoffset+\evensidemargin-\SymbolLegendGap-\SymbolLegendWidth\relax
    \fi\fi
    \SL@y=\dimexpr1in+\voffset+\topmargin+\headheight+\headsep\relax
    \put(\strip@pt\SL@x,-\strip@pt\SL@y){%
      \makebox(0,0)[lt]{\box\SL@box}}%
    \endgroup
  \fi
}
\AtBeginDocument{%
  \ifSymbolLegendEnabled
    % Reserve outer space, preserving the body width when it already fits.
    \SL@available=\dimexpr\paperwidth-\SymbolLegendInnerMargin
      -\SymbolLegendEdgeMargin-\SymbolLegendWidth-\SymbolLegendGap\relax
    \ifdim\SL@available<60mm
      \PackageError{symbol-legend}{Insufficient room for text and legend}%
        {Adjust the page size or the legend dimensions.}%
    \fi
    \ifdim\textwidth>\SL@available\textwidth=\SL@available\fi
    \oddsidemargin=\dimexpr\SymbolLegendInnerMargin-1in-\hoffset\relax
    \evensidemargin=\dimexpr\paperwidth-\SymbolLegendInnerMargin
      -\textwidth-1in-\hoffset\relax
    \hsize=\textwidth\linewidth=\textwidth\columnwidth=\textwidth
    % Foreground runs after the kernel's DiscardShipoutBox decision. Collect
    % the original page before typesetting this hook's legend material.
    \AddToHook{shipout/foreground}[symbol-legend]{%
      \directlua{symbollegend.collect(\number\ShipoutBox)}%
      \SL@render}%
  \fi
}
\makeatother

\ifnum\NullproductsSymbolLegend=1 \SymbolLegendEnabledtrue\fi
\SymbolLegendCompacttrue

\NewDocumentCommand\NPOrd{s m m}{%
  \mathord{\IfBooleanT{#1}{\SymbolLegendIntroduce{#2}}%
    \SymbolLegendUse{#2}#3}}
\NewDocumentCommand\NPRel{s m m}{%
  \mathrel{\IfBooleanT{#1}{\SymbolLegendIntroduce{#2}}%
    \SymbolLegendUse{#2}#3}}
\NewDocumentCommand\NPBin{s m m}{%
  \mathbin{\IfBooleanT{#1}{\SymbolLegendIntroduce{#2}}%
    \SymbolLegendUse{#2}#3}}
\ExplSyntaxOn
\NewDocumentCommand\NPNullIndex{m}{%
  \str_if_eq:nnTF {#1}{Bsp}{\mathrm{Bsp}}{#1}}
\ExplSyntaxOff
\NewDocumentCommand\NPNullPreorder{s m}{%
  \mathrel{\IfBooleanT{#1}{\SymbolLegendIntroduce{nullpreorder-#2}}%
    \SymbolLegendUse{nullpreorder-#2}\preceq_{\NPNullIndex{#2}}}}
\NewDocumentCommand\NPNullEquivalent{s m}{%
  \mathrel{\IfBooleanT{#1}{\SymbolLegendIntroduce{nullequivalence-#2}}%
    \SymbolLegendUse{nullequivalence-#2}\sim_{#2}}}
\NewDocumentCommand\NPExampleS{s}{%
  \IfBooleanTF{#1}{\NPOrd*{example-semigroup}{S_{\mathrm{Bsp}}}}%
    {\NPOrd{example-semigroup}{S_{\mathrm{Bsp}}}}}

\SymbolLegendDeclareReading{example-semigroup}{\(S_{\mathrm{Bsp}}\)}%
  {Dreierhalbgruppe \(\{0,c,a\}\): \(aa=c\); sonst immer \(0\).}
\SymbolLegendDeclareReading{products-example}{\(\mathcal M_{\mathrm{Bsp}}\)}%
  {Die drei Produktmengen \(\{0\},\{c\},\{0,c\}\) im Dreierbeispiel.}
\SymbolLegendDeclareReading{profile-example}{\([A]_{\mathrm{Bsp}}\)}%
  {Profil in \(S_{\mathrm{Bsp}}\): gleiche Nulltests wie \(A\); \(p_0\) oder \(p_1\).}
\SymbolLegendDeclareReading{profiles-example}{\(Q_{\mathrm{Bsp}}\)}%
  {Die zwei Profile \(\{p_0,p_1\}\) des Dreierbeispiels.}
\SymbolLegendDeclareReading{nullpreorder-Bsp}{\(\preceq_{\mathrm{Bsp}}\)}%
  {Profilvergleich im Dreierbeispiel: links und rechts mindestens dieselben Nulltests.}
\SymbolLegendDeclareReading{lower-family-example}{\(L_p^{\mathrm{Bsp}}\)}%
  {Nichtleere Mengen in \(S_{\mathrm{Bsp}}\) mit Profil höchstens \(p\).}
\SymbolLegendDeclareReading{lower-set-example}{\(D_p^{\mathrm{Bsp}}\)}%
  {Grundelemente des Dreierbeispiels mit Singleton-Profil höchstens \(p\).}
\SymbolLegendDeclareReading{fiber-example}{\(C_p^{\mathrm{Bsp}}\)}%
  {Grundelemente des Dreierbeispiels mit Singleton-Profil genau \(p\).}
\SymbolLegendDeclareReading{strict-set-example}{\(E_p^{\mathrm{Bsp}}\)}%
  {Grundelemente mit Profil strikt unter \(p\); also \(D_p^{\mathrm{Bsp}}\setminus C_p^{\mathrm{Bsp}}\).}

\SymbolLegendDeclareReading{phi}{\(\Phi\)}%
  {Ordnet nichtleere Mengen bijektiv zu; \(\Phi(AB)=\Phi(A)\Phi(B)\).}
\SymbolLegendDeclareReading{element-map}{\(f\)}%
  {Gesucht: Bijektion \(S\to T\) mit \(f(xy)=f(x)f(y)\).}
\SymbolLegendDeclareReading{tau}{\(\tau\)}%
  {Vertauscht \(\{a\}\) und \(\{0,a\}\); \(\{0\}\) bleibt fest.}
\SymbolLegendDeclareReading{example-zero}{\(p_0\)}%
  {Im Dreierbeispiel alle nichtleeren Teilmengen von \(\{0,c\}\).}
\SymbolLegendDeclareReading{example-one}{\(p_1\)}%
  {Die vier Teilmengen von \(\{0,c,a\}\), die \(a\) enthalten.}
\SymbolLegendDeclareReading{phi-square}{\(\Phi_2\)}%
  {Wie \(\Phi\), nur für \(U\in\mathcal M_S\): \(\Phi_2(U)=\Phi(U)\).}
\SymbolLegendDeclareReading{zero-target}{\(0_T\)}%
  {Element mit \(0_Tt=t0_T=0_T\) für jedes \(t\in T\).}
\SymbolLegendDeclareReading{nonzero-target}{\(c_T\)}%
  {Zweiter Produktwert: \(T^2=\{0_T,c_T\}\), \(c_T\ne0_T\).}
\SymbolLegendDeclareReading{profile-order}{\(\le\)}%
  {\(p\le q\): Mengen in \(p\) bestehen alle Nulltests der Mengen in \(q\).}
\SymbolLegendDeclareReading{profile-strict}{\(<\)}%
  {\(p<q\) heißt \(p\le q\) und \(p\ne q\).}
\SymbolLegendDeclareReading{profile-map}{\(\psi\)}%
  {\([A]_S\!\mapsto\![\Phi(A)]_T\). Das Profil folgt dem Bild unter \(\Phi\).}
\SymbolLegendDeclareReading{disjoint-union}{\(\mathbin{\dot\cup}\)}%
  {\(A\cup B\), wobei \(A\cap B=\varnothing\): kein Element liegt in beiden.}
\SymbolLegendDeclareReading{lower-map}{\(g_p\)}%
  {Bijektion von \(E_p^S\)\newline nach \(E_{\psi(p)}^T\): alle kleineren Fasern zusammen.}
\SymbolLegendDeclareReading{fiber-map}{\(f_p\)}%
  {Bijektion von \(C_p^S\)\newline nach \(C_{\psi(p)}^T\).}
\SymbolLegendDeclareReading{product-square-S}{\(S^2\)}%
  {Alle Werte \(xy\) mit \(x,y\in S\); hier \(\{0,c\}\).}
\SymbolLegendDeclareReading{product-square-T}{\(T^2\)}%
  {Alle Werte \(xy\) mit \(x,y\in T\); hier \(\{0_T,c_T\}\).}
\SymbolLegendDeclareReading{product-square-N}{\(N^2\)}%
  {Alle Werte \(xy\) mit \(x,y\in N\); im Beispiel nur \(0\).}
\SymbolLegendDeclareReading{products-S}{\(\mathcal M_S\)}%
  {Alle Mengen \(AB\) mit nichtleeren \(A,B\subseteq S\).}
\SymbolLegendDeclareReading{nullpreorder-S}{\(\preceq_S\)}%
  {Aus \(BX=\{0\}\) folgt \(AX=\{0\}\), für jeden Test \(X\); links ebenso.}
\SymbolLegendDeclareReading{nullequivalence-S}{\(\sim_S\)}%
  {\(A,B\) geben bei jedem linken und rechten Nulltest dieselbe Antwort.}
\SymbolLegendDeclareReading{profile-S}{\([A]_S\)}%
  {Alle nichtleeren \(B\subseteq S\), die dieselben Nulltests wie \(A\) bestehen.}
\SymbolLegendDeclareReading{profiles-S}{\(Q_S\)}%
  {Alle verschiedenen Klassen \([A]_S\) für nichtleere \(A\subseteq S\).}
\SymbolLegendDeclareReading{lower-family-S}{\(L_p^S\)}%
  {Alle nichtleeren \(A\subseteq S\) mit \([A]_S\le p\).}
\SymbolLegendDeclareReading{lower-set-S}{\(D_p^S\)}%
  {Alle \(s\in S\) mit\newline \([\{s\}]_S\le p\).}
\SymbolLegendDeclareReading{fiber-S}{\(C_p^S\)}%
  {Alle \(s\in S\) mit\newline \([\{s\}]_S=p\).}
\SymbolLegendDeclareReading{strict-set-S}{\(E_p^S\)}%
  {Alle \(s\in S\) mit\newline \([\{s\}]_S<p\); also \(D_p^S\setminus C_p^S\).}
\SymbolLegendDeclareReading{products-T}{\(\mathcal M_T\)}%
  {Alle Mengen \(AB\) mit nichtleeren \(A,B\subseteq T\).}
\SymbolLegendDeclareReading{nullpreorder-T}{\(\preceq_T\)}%
  {Aus \(BX=\{0_T\}\) folgt \(AX=\{0_T\}\), für jeden Test \(X\); links ebenso.}
\SymbolLegendDeclareReading{profile-T}{\([A]_T\)}%
  {Alle nichtleeren \(B\subseteq T\), die dieselben Nulltests wie \(A\) bestehen.}
\SymbolLegendDeclareReading{profiles-T}{\(Q_T\)}%
  {Alle verschiedenen Klassen \([A]_T\) für nichtleere \(A\subseteq T\).}
\SymbolLegendDeclareReading{lower-family-T}{\(L_r^T\)}%
  {Alle nichtleeren \(A\subseteq T\) mit \([A]_T\le r\).}
\SymbolLegendDeclareReading{lower-set-T}{\(D_r^T\)}%
  {Alle \(t\in T\) mit\newline \([\{t\}]_T\le r\).}
\SymbolLegendDeclareReading{fiber-T}{\(C_r^T\)}%
  {Alle \(t\in T\) mit\newline \([\{t\}]_T=r\).}
\SymbolLegendDeclareReading{strict-set-T}{\(E_r^T\)}%
  {Alle \(t\in T\) mit\newline \([\{t\}]_T<r\); also \(D_r^T\setminus C_r^T\).}
\SymbolLegendDeclareReading{ideal-L}{\(L_H(u)\)}%
  {\(u\) und alle Produkte \(xu\) mit \(x\in H\).}
\SymbolLegendDeclareReading{ideal-R}{\(R_H(u)\)}%
  {\(u\) und alle Produkte \(uy\) mit \(y\in H\).}
\SymbolLegendDeclareReading{ideal-J}{\(J_H(u)\)}%
  {\(u\) und alle Produkte \(xu,uy,xuy\) mit \(x,y\in H\).}
\SymbolLegendDeclareReading{green-L}{\(\mathcal L_H\)}%
  {\(u,v\) sind verwandt, wenn \(L_H(u)=L_H(v)\).}
\SymbolLegendDeclareReading{green-R}{\(\mathcal R_H\)}%
  {\(u,v\) sind verwandt, wenn \(R_H(u)=R_H(v)\).}
\SymbolLegendDeclareReading{green-J}{\(\mathcal J_H\)}%
  {\(u,v\) sind verwandt, wenn \(J_H(u)=J_H(v)\).}
\SymbolLegendDeclareReading{green-H}{\(\mathcal H_H\)}%
  {Es gelten zugleich \(L_H(u)=L_H(v)\) und \(R_H(u)=R_H(v)\).}
\SymbolLegendDeclareReading{ann-right}{\(\operatorname{Ann}_r(A)\)}%
  {Alle \(x\in S\) mit \(ax=0\) für jedes \(a\in A\).}
\SymbolLegendDeclareReading{ann-left}{\(\operatorname{Ann}_l(A)\)}%
  {Alle \(x\in S\) mit \(xa=0\) für jedes \(a\in A\).}
\SymbolLegendDeclareReading{example-dot}{\(\cdot\)}%
  {Im Fünferbeispiel: \(aa=bb=u\), \(ab=ba=v\); sonst \(0\).}
\SymbolLegendDeclareReading{example-star}{\(\ast\)}%
  {Im Fünferbeispiel: \(aa=ab=bb=u\), \(ba=v\); sonst \(0\).}
